\documentclass[10pt,a4paper]{article}

\usepackage[margin=17mm]{geometry}
\usepackage[T1]{fontenc}
\usepackage{lmodern}
\usepackage{amsmath,amssymb}
\usepackage{graphicx}
\usepackage{booktabs,tabularx,array}
\usepackage{xcolor}
\usepackage{fancyhdr}
\usepackage{titlesec}
\usepackage{microtype}
\usepackage{hyperref}

\definecolor{navy}{HTML}{174E74}
\definecolor{muted}{HTML}{586875}

\hypersetup{
  colorlinks=true,
  linkcolor=navy,
  urlcolor=navy,
  pdftitle={Discrete OPC Oscillation Theory:
    Recurrence, Modes and Scaling}
}

\renewcommand{\familydefault}{\sfdefault}
\setlength{\parindent}{0pt}
\setlength{\parskip}{6pt}
\setlength{\headheight}{14pt}
\setlength{\emergencystretch}{2em}
\renewcommand{\arraystretch}{1.22}

\pagestyle{fancy}
\fancyhf{}
\fancyhead[L]{\small\color{muted}
  DISCRETE OPC / TECHNICAL NOTE}
\fancyhead[R]{\small\color{muted}
  23 SEPTEMBER 2026}
\fancyfoot[L]{\footnotesize\color{muted}
  Final mask correction under a linear, finite-range model}
\fancyfoot[R]{\footnotesize\color{muted}\thepage}
\renewcommand{\headrulewidth}{0pt}
\renewcommand{\footrulewidth}{0.3pt}

\titleformat{\section}
  {\Large\bfseries\color{navy}}
  {\thesection.}{0.5em}{}
\titleformat{\subsection}
  {\normalsize\bfseries\color{navy}}
  {}{0pt}{}

\titlespacing*{\section}{0pt}{0pt}{10pt}
\titlespacing*{\subsection}{0pt}{10pt}{4pt}

\newcolumntype{Y}{>{\raggedright\arraybackslash}X}
\newcommand{\subtitle}[1]{%
  {\small\color{muted}#1\par}\medskip}
\newcommand{\smallnote}[1]{%
  {\small #1\par}}

\begin{document}

% ============================================================
\section{Model and correction recurrence}
\subtitle{A discrete explanation for damped spatial
oscillations in a converged mask CD profile.}

The object of interest is the \textbf{final mask correction},
not a residual wafer-CD oscillation.
A localized error at the left block edge can excite damped
spatial modes of the inverse model.
We derive their amplitude, period and decay without
prescribing a sinusoid as the solution.

\begin{tabularx}{\textwidth}{@{}p{0.29\textwidth}Y@{}}
\toprule
\textbf{Symbol} & \textbf{Meaning} \\
\midrule
$j=0,1,\ldots$; $s=jd$
&
Control index and physical position;
$j=0$ is the first control.
\\
$u_j^{(n)},\,e_j^{(n)}$
&
Accumulated mask-CD correction and wafer-CD error
at iteration $n$.
\\
$d$; $M$; $N=2M+1$
&
Control spacing; number of neighbors on either side;
full stencil size.
\\
$R=Md$; $R_c=Nd/2$
&
Range to the outer control center;
cell-based half-width $R+d/2$.
\\
$a>0$; $b>0$; $E$
&
Local sensitivity; total nonlocal sensitivity;
initial error at $j=0$.
\\
\bottomrule
\end{tabularx}

The nonlocal kernel is uniform and normalized.
Set controls outside the left boundary to zero;
do not renormalize the shortened stencil there.
All interior controls have the same sensitivities.

\begin{equation}
e_j^{(n)}
=
E\delta_{j0}
+
a u_j^{(n)}
+
\frac{b}{N}
\sum_{\ell=-M}^{M}u_{j+\ell}^{(n)}.
\end{equation}

\begin{equation}
u_j^{(n+1)}
=
u_j^{(n)}-\eta e_j^{(n)},
\qquad
u_j^{(0)}=0,
\qquad
u_j^{(n)}=0\quad(j<0).
\end{equation}

The Kronecker delta is $1$ at $j=0$ and $0$ elsewhere.
The update is \textbf{P-type error correction}.
It is not gradient descent on squared error, which would
additionally apply the transposed sensitivity matrix.

\begin{equation}
\mathbf e^{(n)}
=
E\mathbf e_0+J\mathbf u^{(n)},
\qquad
\mathbf e^{(n+1)}
=
(I-\eta J)\mathbf e^{(n)}.
\end{equation}

\smallnote{
Here the bold vector $\mathbf e_0=(1,0,\ldots)$ selects
the boundary control.
For a finite positive-definite $J$, convergence requires
$0<\eta<2/\lambda_{\max}(J)$.
The experiments used $\eta=1/(a+b)$.
Spatial decay below is distinct from decay with
iteration count.
}

% ============================================================
\clearpage
\section{Fixed point and spatial modes}
\subtitle{Why powers of a characteristic root span
the solution.}

If the correction converges and $\eta$ is nonzero,
successive corrections become equal and wafer error
tends to zero.
The equilibrium equation is therefore

\begin{equation}
a u_j
+
\frac{b}{N}
\sum_{\ell=-M}^{M}u_{j+\ell}
=
-E\delta_{j0},
\qquad j\geq 0.
\end{equation}

For $j\geq M$ the stencil is complete and the
right-hand side is zero.
Solving for the rightmost value gives a spatial
recurrence of order $2M$:

\begin{equation}
u_{j+M}
=
-\frac{Na}{b}u_j
-
\sum_{\ell=-M}^{M-1}u_{j+\ell},
\qquad j\geq M.
\end{equation}

Thus the first $2M$ values uniquely determine
all subsequent values.
To find a basis for this recurrence, insert a
geometric sequence $u_j=r^j$ and divide by $r^j$.
Multiplication by $r^M$ removes negative powers:

\begin{equation}
a+\frac{b}{N}\sum_{\ell=-M}^{M}r^\ell=0
\quad\Longrightarrow\quad
P_M(r)
=
a r^M+\frac{b}{N}\sum_{h=0}^{2M}r^h
=
0.
\end{equation}

Each root supplies a solution $r^j$.
If the $2M$ roots are distinct, their first $2M$
values form a Vandermonde matrix.
Its determinant is nonzero:

\begin{equation}
V_{k\nu}=r_\nu^k
\quad(k=0,\ldots,2M-1),
\qquad
\det V
=
\prod_{\mu<\nu}(r_\nu-r_\mu)
\neq 0.
\end{equation}

Consequently any initial $2M$ values can be expressed
uniquely by these modes.
Both the resulting mode sum and the desired sequence
satisfy the same recurrence and initial values,
so they agree at every later index.
This establishes completeness, not merely an
exponential trial solution.

\begin{equation}
u_j=\sum_{\nu=1}^{2M}C_\nu r_\nu^j.
\end{equation}

The polynomial is reciprocal:
its coefficients read the same in reverse.
Roots occur as $r$ and $1/r$, and real coefficients
also imply conjugate roots.
In the stable regime with no unit-circle roots,
$M$ roots lie inside the unit circle.
Requiring $u_j$ to vanish as $j$ tends to infinity
eliminates the growing modes:

\begin{equation}
u_j=\sum_{\nu=1}^{M}C_\nu r_\nu^j,
\qquad |r_\nu|<1.
\end{equation}

\smallnote{
This note assumes distinct roots, as in the numerical
examples.
A repeated root requires additional terms $j^p r^j$.
The semi-infinite derivation describes a finite block
when the far boundary has negligible influence
in the observation window.
}

% ============================================================
\clearpage
\section{Period and spatial decay}
\subtitle{Exact relations from a conjugate root pair;
equivalent $q$ representation.}

The real boundary problem has a unique decaying solution
in the stable regime.
Complex conjugation produces the same solution,
so coefficients belonging to conjugate roots
are conjugates.
Select an excited complex root $r_*$ with the largest
modulus below $1$, and choose its positive argument.

\begin{equation}
r_*=\rho e^{i\theta},
\qquad
C_*=|C_*|e^{i\phi},
\qquad
0<\rho<1,
\qquad
0<\theta<\pi.
\end{equation}

\begin{equation}
C_*r_*^j
+
\overline{C_*}\,\overline{r_*}^{\,j}
=
2|C_*|\rho^j\cos(j\theta+\phi).
\end{equation}

With $s=jd$, the envelope is
$\exp[(s/d)\ln\rho]$ and the phase advances by
$\theta/d$ per unit length.
A full period requires a phase change of $2\pi$.
Therefore

\begin{equation}
\boxed{
A_0=2|C_*|,
\qquad
L=\frac{2\pi d}{\theta},
\qquad
\gamma=-\frac{\ln\rho}{d},
\qquad
\phi=\arg C_*.
}
\end{equation}

\begin{equation}
u_{\mathrm{dom}}(s)
=
A_0 e^{-\gamma s}
\cos\left(\frac{2\pi s}{L}+\phi\right).
\end{equation}

This is the dominant pair extracted from the exact
mode sum.
Other roots contribute faster-decaying terms and may
matter near the boundary.
$A_0$ is the dominant envelope extrapolated to $s=0$;
it is not the actual value $u_0$.

\subsection*{From $r$ to the dimensionless complex frequency $q$}

\begin{equation}
r=e^{2iq/N},
\qquad
q=\frac{N}{2}\bigl(\theta-i\ln\rho\bigr).
\end{equation}

The finite geometric sum can be evaluated
without approximation:

\begin{equation}
\sum_{\ell=-M}^{M}e^{2i\ell q/N}
=
e^{-2iMq/N}
\frac{1-e^{2iq}}{1-e^{2iq/N}}
=
\frac{\sin q}{\sin(q/N)}.
\end{equation}

Using $1-e^{2ix}=-2i e^{ix}\sin x$,
the exponential prefactors cancel because $N=2M+1$.
Hence, with $c=a/b$,

\begin{equation}
c+F_N(q)=0,
\qquad
F_N(q)=\frac{\sin q}{N\sin(q/N)},
\qquad
q_*=\alpha_N+i\beta_N.
\end{equation}

\begin{equation}
\boxed{
L=\frac{\pi Nd}{\alpha_N}
 =\frac{2\pi R_c}{\alpha_N},
\qquad
\gamma=\frac{2\beta_N}{Nd}
      =\frac{\beta_N}{R_c}.
}
\end{equation}

\smallnote{
The roots $\alpha_N$ and $\beta_N$ depend on both $c$
and $N$.
Replacing $N\sin(q/N)$ by $q$ would be a continuum
approximation; it is not used for the exact
discrete predictions.
}

% ============================================================
\clearpage
\section{Amplitude from the left boundary}
\subtitle{The source fixes the mode coefficients;
a characteristic root alone does not.}

For $j=0,\ldots,M-1$, the left part of the stencil
lies outside the domain.
Since those controls are zero, the first $M$
equilibrium equations are

\begin{equation}
a u_j
+
\frac{b}{N}\sum_{k=0}^{j+M}u_k
=
-E\delta_{j0},
\qquad
j=0,\ldots,M-1.
\end{equation}

Substitute the complete decaying mode sum and
interchange the finite sums:

\begin{equation}
\sum_{\nu=1}^{M}
\left[
a r_\nu^j
+
\frac{b}{N}\sum_{k=0}^{j+M}r_\nu^k
\right]C_\nu
=
-E\delta_{j0}.
\end{equation}

Define an $M\times M$ boundary matrix $B$,
with row index $j=0,\ldots,M-1$ and
column index $\nu=1,\ldots,M$:

\begin{equation}
B_{j\nu}
=
a r_\nu^j
+
\frac{b}{N}
\frac{1-r_\nu^{j+M+1}}{1-r_\nu}.
\end{equation}

\begin{equation}
B\mathbf C=-E\mathbf e_0,
\qquad
\mathbf C=-E B^{-1}\mathbf e_0.
\end{equation}

The geometric-sum expression is valid since a
decaying root is not $1$.
All interior equations are already satisfied
by the roots, and these $M$ equations fix the
$M$ remaining coefficients.
Numerically, solve the linear system rather than
forming an inverse.

\begin{equation}
A_0
=
2|E|
\left|
\left(B^{-1}\mathbf e_0\right)_*
\right|.
\end{equation}

\subsection*{Separate the error scale from the model parameters}

\begin{equation}
B=b\widetilde B,
\qquad
\widetilde B_{j\nu}
=
c r_\nu^j
+
\frac{1}{N}
\frac{1-r_\nu^{j+M+1}}{1-r_\nu}.
\end{equation}

\begin{equation}
\boxed{
A_0=\frac{|E|}{b}\,\mathcal G(c,M),
\qquad
\mathcal G(c,M)
=
2\left|
\left(\widetilde B^{-1}\mathbf e_0\right)_*
\right|.
}
\end{equation}

\textbf{Exact consequence:}
For a fixed model, multiplying $E$ by a factor
multiplies the entire correction profile by
that factor.
The period and spatial decay do not change;
envelope and peak-to-peak magnitudes scale by
its absolute value.

\smallnote{
\textbf{Amplitude convention:}
The experimentally discussed second-largest cycle
peak-to-peak is a separate metric.
Near the block edge, several modes overlap,
so it cannot generally be replaced by $2A_0$.
Even a single damped mode changes its envelope
between successive extrema.
}

% ============================================================
\clearpage
\section{Critical point and finite-range decay}
\subtitle{Distinguish oscillation, spatial damping and
convergence of the correction loop.}

A positive real $r$ between $0$ and $1$ produces
monotone decay; a negative real $r$ between $-1$
and $0$ alternates every sample; a nonreal $r$
produces an oscillation with period
$2\pi d/\arg(r)$.
For $a,b>0$, $P_M(r)$ has only positive terms
when $r>0$, so this box model has no positive-real
spatial root.

Define the real-frequency model response and
its bulk stability threshold:

\begin{equation}
H(\theta)
=
a+\frac{b}{N}
\left[
1+2\sum_{\ell=1}^{M}\cos(\ell\theta)
\right].
\end{equation}

\begin{equation}
c_N
=
-\min_{0\leq\theta\leq\pi}
\frac{1}{N}
\left[
1+2\sum_{\ell=1}^{M}\cos(\ell\theta)
\right],
\qquad
H_{\min}=b(c-c_N).
\end{equation}

In the tested stencils, the deepest negative lobe
is the first one.
Let $q_N$ be its minimum in the equivalent
$F_N(q)$ representation.
Differentiation gives

\begin{equation}
F_N(q_N)=-c_N,
\qquad
\cot q_N=\frac{1}{N}\cot(q_N/N).
\end{equation}

\begin{equation}
\frac{F_N'}{F_N}
=
\cot q-\frac{1}{N}\cot(q/N),
\qquad
F_N''(q_N)=c_N(1-N^{-2}).
\end{equation}

For the last equality, differentiate the logarithmic
derivative.
At the minimum its square vanishes, and the
cotangent identity cancels the remaining
cotangent-squared terms.
A Taylor expansion then yields

\begin{equation}
0
\simeq
(c-c_N)
+
\frac{c_N(1-N^{-2})}{2}(q-q_N)^2.
\end{equation}

\begin{equation}
\boxed{
L\simeq\frac{2\pi R_c}{q_N},
\qquad
\gamma\simeq
\frac{1}{R_c}
\sqrt{
\frac{2(c-c_N)}
     {c_N(1-N^{-2})}
}.
}
\end{equation}

This approximation applies just above $c_N$.
As $c$ approaches $c_N$ from above, the pair
approaches the real $q$ axis:
spatial damping tends to zero while the period
stays finite.
Real $q$ means $|r|=1$, not a real $r$.
Below the threshold, the nearby real-frequency
modes are undamped and some $H(\theta)$
become negative.

\smallnote{
For P correction, a bulk error mode is multiplied
by $1-\eta H(\theta)$.
Negative $H$ therefore causes growth for positive
$\eta$.
The threshold concerns bulk stability and loss
of damping, not the first appearance of any
oscillation.
A finite block has its own exact eigenvalue threshold.
}

\smallnote{
\textbf{Why decay can be nonmonotonic in $R$:}
Increasing $M$ decreases $c_N$ toward its continuum
value.
Near the threshold this increases the square-root
factor strongly, while $1/R_c$ decreases.
Their competition produces the rise and fall
seen in the numerical sweep.
}

% ============================================================
\clearpage
\section{What is proportional to what?}
\subtitle{Exact discrete dependencies first;
large-ambit approximations second.}

\begin{equation}
L=d\,\mathcal L(c,M),
\qquad
\gamma=\frac{1}{d}\,g(c,M),
\qquad
A_0=\frac{|E|}{b}\,\mathcal G(c,M).
\end{equation}

\begin{equation}
\mathcal L(c,M)
=
\frac{2\pi}{\arg r_*},
\qquad
g(c,M)=-\ln|r_*|,
\qquad
c=a/b.
\end{equation}

{\hbadness=10000
\small
\begin{tabularx}{\textwidth}{@{}
  p{0.36\textwidth}YYY@{}}
\toprule
\textbf{Change / fixed quantities}
&
\textbf{Amplitude $A_0$}
&
\textbf{Period $L$}
&
\textbf{Decay $\gamma$}
\\
\midrule
Change $|E|$; $a,b,M,d$ fixed
&
Proportional to $|E|$
&
Unchanged
&
Unchanged
\\
\addlinespace
Scale $a,b$ by $k>0$; $E,M,d$ fixed
&
Divided by $k$
&
Unchanged
&
Unchanged
\\
\addlinespace
Change $d$; $a,b,M,E$ fixed
&
Unchanged
&
Proportional to $d$
&
Proportional to $1/d$
\\
\addlinespace
Change $M$ (ambit); $a,b,d,E$ fixed
&
$\mathcal G(c,M)$ determines it
&
Root angle determines it
&
Root modulus determines it
\\
\addlinespace
Change $a/b$ at fixed $b,M,d,E$
&
No universal power law
&
No universal power law
&
Strong near-threshold sensitivity
\\
\bottomrule
\end{tabularx}\par
}

The third row rescales physical spacing while
keeping the same discrete stencil.
It is different from extending the ambit by
adding neighbors at fixed $d$.
In the fourth row, no exact universal power law
in $R$ is implied.

\subsection*{Large ambit: fixed $d$ and fixed normalized
coupling strengths}

When $N$ is large, $F_N(q)$ approaches $\sin(q)/q$.
At fixed $c$ above the continuum threshold,
the dominant $q$ approaches a fixed complex value
$\alpha_\infty+i\beta_\infty$.
Thus

\begin{equation}
L\simeq\frac{2\pi}{\alpha_\infty}R_c,
\qquad
\gamma\simeq\frac{\beta_\infty}{R_c},
\qquad
R_c=R+d/2.
\end{equation}

For the same one-cell source, its physical width
$d$ and integrated strength $Ed$ remain fixed.
In the narrow-source continuum limit of this
half-line boundary problem, the dimensionless
amplitude gain is proportional to $d/R_c$:

\begin{equation}
A_0
\simeq
K(c)\frac{|E|d}{bR_c},
\qquad
d/R_c\ll 1.
\end{equation}

$K(c)$ is fixed by the normalized kernel and
the boundary problem.
This amplitude relation is an asymptotic
boundary-response result, not an exact identity
at finite $M$.
The exact discrete calculation remains the
boundary-matrix solve in Section~4.

\begin{equation}
c=0.22:
\qquad
L\simeq1.398897R_c,
\qquad
\gamma\simeq0.159433/R_c.
\end{equation}

\smallnote{
\textbf{Interpretation:}
At large ambit the extrapolated envelope decreases
approximately as $1/R_c$, but the decay length
$1/\gamma$ grows with $R_c$.
At a fixed distant position,
$A_0\exp(-\gamma s)$ can therefore increase even
when $A_0$ decreases.
}

% ============================================================
\clearpage
\section{Numerical validation and use}
\subtitle{Independent correction experiments versus
discrete characteristic-root predictions.}

The sweep used $a=0.22$, $b=1$,
$d=0.08\,\mu\mathrm{m}$, $E=1\,\mathrm{nm}$,
and $M=8,9,10,12,14,18,24$.
Each experiment started from $u=0$ and iterated
the P update to
$\max|e|/|E|\leq10^{-12}$ on 4096 controls.
There was no periodic wrap.

\smallnote{
Measurements came from a dominant-mode fit over
$s=8$--$32\,\mu\mathrm{m}$, initialized from
measured peaks rather than theoretical roots.
Predictions used all decaying roots and the
boundary linear system.
The three curves agree without fitting theory
parameters to the profiles.
}

\begin{center}
\includegraphics[
  width=0.91\textwidth,
  height=0.45\textheight,
  keepaspectratio
]{validation_portrait.png}
\end{center}

{\hbadness=10000
\small
\begin{tabularx}{\textwidth}{@{}YYY@{}}
\toprule
\textbf{Quantity}
&
\textbf{Maximum relative difference}
&
\textbf{Baseline theory: $M=9$}
\\
\midrule
Envelope amplitude $A_0$
&
$0.008821\%$
&
$1.35087026\,E$
\\
Spatial decay $\gamma$
&
$0.005320\%$
&
$0.10720785\,\mu\mathrm{m}^{-1}$
\\
Period $L$
&
$0.000791\%$
&
$1.06184467\,\mu\mathrm{m}$
\\
\bottomrule
\end{tabularx}\par
}

\smallnote{
Small remaining differences include finite-window
extraction of a dominant mode from a multimode
profile.
An independent banded solve also verified
the converged iterations;
the largest absolute difference was
$4.4\times10^{-9}\,\mathrm{nm}$.
}

\smallnote{
\textbf{Practical recipe:}
Specify $a,b,d,M$ and the boundary error;
solve the characteristic polynomial;
retain decaying roots;
solve the boundary system;
select the slowest excited pair;
compute $A_0$, $L$ and $\gamma$.
This validates the mechanism for the stated
linear box model.
It demonstrates that converged mask correction
need not be spatially flat;
it does not by itself identify the kernel of
a production lithography model.
}

\end{document}
